\chapter{Calibration}
\label{ch:calibration}
\index{calibration}

\section{Overview}
\label{sec:cal-overview}

Stimulus levels in pyabr3 depend on a speaker calibration file. The file holds, for each
test frequency, the sound level that a full-scale (10~V) tone produces with no
attenuation. Making this file is a three-stage procedure, with an optional fourth
measurement:

\begin{enumerate}
  \item \textbf{Microphone:} measure the microphone's sensitivity using a sound level
        calibrator.
  \item \textbf{Speaker:} measure the sound level at each frequency at a fixed
        attenuator setting.
  \item \textbf{75~dB check:} use the new calibration to request 75~dB~SPL at every
        frequency, and measure what is produced.
  \item \textbf{Noise floor} (optional): record the microphone with no stimulus.
\end{enumerate}

\needsinput{How often the lab calibrates, and the criteria for accepting a calibration
(for example, how far the 75~dB check may deviate from 75~dB, and at which
frequencies).}

\begin{caution}
\textbf{Validation against MATLAB is not finished.}\index{validation}
The calibration code is a port of the ABR4 MATLAB program. Its formulas reproduce
stored MATLAB results, and it has been tested on simulated hardware. It has not yet
been run on the rig or compared with MATLAB on the same recordings. Until that
comparison has been made, compare new calibrations with recent MATLAB calibrations.
\end{caution}

\section{The calibration window}
\index{calibration window}

Press \gui{Calibrate} in the main window. During an acquisition, a message asks you to
stop it first. The window title shows which hardware is in use.

\begin{caution}
If the RP2.1 is not connected, the window uses \textbf{simulated hardware} and the title
says MOCK. Simulated results are written to a temporary folder, so they cannot replace a
real calibration.
\end{caution}

\figorplaceholder{cal_speaker}{The calibration window after a speaker calibration. This
example used the simulated hardware.}{Calibration window after a speaker run.}

The panel on the left holds the settings:

\begin{center}
\begin{tabularx}{\textwidth}{@{}lX@{}}
\toprule
Setting & Meaning \\
\midrule
\gui{Microphone} & The microphone in use (Table~\ref{tab:microphones}) \\
\gui{Speaker} & The speaker in use. Choosing one fills in the next two settings and
  \gui{Attn} (Table~\ref{tab:speakers}) \\
\gui{Frequencies (kHz)} & Test frequencies, separated by commas or spaces \\
\gui{Ref SPL (maxtones)} & Stored in the file \\
\gui{Cal directory} & Where files are read and written. Defaults to the folder of
  \key{calfile} in the configuration \\
\gui{Test calc (no save)} & For microphone runs: use a synthetic signal, and do not
  save \\
\gui{Mic preamp gain (dB)} & The gain set by hand on the microphone preamplifier \\
\gui{Speaker driver attn (dB)} & The setting on the speaker driver \\
\gui{Attn (dB)} & The programmable attenuator setting for speaker runs \\
\bottomrule
\end{tabularx}
\end{center}

\begin{table}[htbp]
\centering
\caption{Speaker settings. The other names in the list use the settings of the same
family.}
\label{tab:speakers}
\begin{tabular}{@{}llll@{}}
\toprule
Speaker & Frequencies & Attn & Ref SPL \\
\midrule
MF1 family & 41, from 1 to 87~kHz & 20~dB & 110 \\
EC1, ES1 family & 28, from 1 to 64~kHz & 0~dB & 83.9 \\
\bottomrule
\end{tabular}
\end{table}

\begin{table}[htbp]
\centering
\caption{Microphones, with the manufacturer's sensitivity. This is printed in the log so
that it can be compared with each measurement.}
\label{tab:microphones}
\begin{tabular}{@{}lll@{}}
\toprule
Microphone & Sensitivity (mV/Pa) & dB re 1 V/Pa \\
\midrule
7012\#39279 (1/2 inch) & 16.22 & $-35.8$ \\
7016\#9945 (1/4 inch) & 3.55 & $-49.0$ \\
7016\#10252 (1/4 inch) & 3.85 & $-48.3$ \\
\bottomrule
\end{tabular}
\end{table}

Before each run, a dialog lists the settings that the program cannot read from the
hardware (Figure~\ref{fig:cal_verify_dialog}). Check them at the equipment, and answer
\gui{Yes} only if they are correct. The values shown come from the panel on the left.

\figorplaceholder[0.45\textwidth]{cal_verify_dialog}{The settings check shown before a
speaker calibration.}{Settings confirmation dialog.}

\section{Calibration procedure}
\label{sec:cal-procedure}

\subsection{Microphone}
\index{microphone calibration}\index{calibrator}

\needsinput{How the microphone is seated in the calibrator, and the calibrator model.}

\begin{enumerate}
  \item Place the microphone in the calibrator. Set the calibrator to 94~dB and turn it
        on.
  \item Select the \gui{Microphone}, and check \gui{Mic preamp gain (dB)}.
  \item Press \gui{Microphone (94 dB)}, and confirm the settings.
  \item Read the \gui{Log}. Compare the measured transfer factor with the manufacturer's
        value printed above it.
\end{enumerate}

The program records 1~s with the attenuator at maximum. It finds the sensitivity from
the signal in a 940--1020~Hz band. The result is saved as
\file{microphone\_<name>.cal}.

\begin{caution}
\textbf{Use the 94~dB button.}\index{104 dB calibrator} After a calibration with
\gui{Microphone (104 dB)}, speaker levels read 10~dB too high. The log prints a warning
when this applies. This matches the MATLAB program, and will be corrected once
validation is complete.
\end{caution}

\figorplaceholder{cal_microphone}{A microphone calibration, using the simulated
hardware.}{Calibration window after a microphone run.}

\subsection{Speaker}
\index{speaker calibration}

\needsinput{Speaker and microphone placement for calibration, including the distance
between them and how this relates to the position of the animal.}

\begin{enumerate}
  \item Place the microphone at the position of the animal's ear.
  \item Select the \gui{Speaker}. Check \gui{Mic preamp gain (dB)},
        \gui{Speaker driver attn (dB)} and \gui{Attn (dB)}.
  \item Press \gui{Calibrate Speaker}, and confirm the settings.
\end{enumerate}

The run takes at least one minute for the MF1 frequency list. The four plots update
after each frequency. \gui{Stop} ends the run, and nothing is saved.

The run is refused if the preamplifier gain differs from the gain in the microphone's
calibration file. Measure the microphone again at the new gain, or restore the gain.

At each frequency the program:

\begin{enumerate}
  \item plays a 0.5~s, 10~V tone with 5~ms rise and fall times, and records 0.75~s from
        the microphone at 195~kHz;
  \item repeats the recording with the attenuator at maximum, to measure the noise
        floor;
  \item removes power supply noise with notch filters at 19.5--20.5~kHz and
        39--41~kHz\index{notch filters};
  \item measures the level in the middle half of the tone (25\% to 75\% of the tone
        duration), with a band-pass filter of $\pm$1/5 octave and with a cosinor fit.
\end{enumerate}

The level stored for use in stimuli, \key{maxdB}, is the band-pass measurement plus the
attenuator setting. It estimates the level at 0~dB attenuation. A frequency inside a
notch filter band produces a warning in the log, because the tone itself would be
removed.

\subsection{Check at 75 dB}
\index{75 dB check}

Leave the microphone and speaker in place, and press \gui{Check Speaker (75 dB)}. For
each frequency the program sets the attenuator to give 75~dB~SPL according to the new
calibration, then measures the result. A good calibration gives a flat line at 75~dB.
Frequencies where the speaker cannot reach 75~dB use 0~dB attenuation, and fall below
the line.

\figorplaceholder{cal_check}{A 75~dB check, using the simulated hardware.}{Calibration
window after a 75 dB check.}

\subsection{Noise floor}
\index{noise floor}

\gui{Noise Floor} records 15~s from the microphone with the attenuator at maximum. The
recording is saved as a text file, \file{Noise\_Measure\_<date and time>.csv}.

\section{Choosing the calibration to use}
\label{sec:cal-select}
\index{Read Cal File}

A new calibration is not used automatically. Press \gui{Read Cal File} in the main
window, and select a file (Figure~\ref{fig:cal_chooser}). The list shows the current
files and those in the \file{calibration\_history} folder, most recent first.
\gui{Browse for File...} finds files in other folders.

\figorplaceholder[0.6\textwidth]{cal_chooser}{The calibration file chooser.}{Calibration
file chooser.}

The chosen calibration is used for later stimuli and stored in the data files. If the
file cannot be read, a message appears and the current calibration remains in use.
Files in the older formats cannot be read.

The \gui{Calibration} tab shows the selected calibration
(Figure~\ref{fig:calibration_tab}). The upper plot shows the levels at the calibration
attenuation, the estimated maximum level, and the noise floor. The lower plot shows the
75~dB check. The check displayed is the first one made after the calibration on the same
day. If there is none, the title of the lower plot says so.

\figorplaceholder{calibration_tab}{The \gui{Calibration} tab, showing a stored
calibration from December 2024 and its 75~dB check.}{Calibration tab of the main
window.}

\section{Calibration files}
\index{calibration files}

All of these are written to the \gui{Cal directory}. Each \file{.cal} file is a MATLAB
file with the same layout as those written by ABR4. Each time one is written, a copy is
placed in \file{calibration\_history/}, with the date and time in its name.

\begin{center}
\begin{tabularx}{\textwidth}{@{}lX@{}}
\toprule
File & Contents \\
\midrule
\file{microphone\_<name>.cal} & Microphone sensitivity \\
\file{frequency\_<speaker>.cal} & Speaker calibration \\
\file{chk75db\_<speaker>.cal} & 75~dB check \\
\file{Noise\_Measure\_<date and time>.csv} & Noise floor recording (no history copy) \\
\bottomrule
\end{tabularx}
\end{center}
